\subsection{The trace-zero column and its slope}
\label{sec:ramified-column}

The level-zero metric in \eqref{eq:lower-slopes} controls the full column. The division argument uses the kernel of the level trace, so its slope also involves the index of the trace-zero lattice. Use \(K\) for the coefficient field in
this calculation. The definite unitary group has rank \(k\ge2\) over a
totally real field of degree \(f\); its CM field has degree \(2f\).
The fine level at each of the two auxiliary type primes \(g_i\)
is the inverse image of an upper unitriangular subgroup. The prime assigned to the sequence splits completely, giving \(f\) places in the totally real field.

Fix a word \(u\) and put \(T=T_u\). Write \(S\) for the level-zero
column lattice in its coefficient space over \(K\), and put
\(d_S=d_u^S\).

\begin{proposition}[Rank and slope of the trace-zero column]\label{prop:trace-zero-column}
Let \(S\) be the level-zero column lattice of the selected
word \(u\), with the fine level at each \(g_i\) the inverse image
of the upper unitriangular subgroup. Suppose the assigned
prime \(T=T_u\ge3\) splits completely in the CM field. Impose the trace-zero condition at each of the \(f\) primes of its maximal totally real subfield above \(T\). The resulting lattice \(S_R\), with the
induced metric, satisfies
\[
 \operatorname{rk} S_R=d_S(1-1/(T-1))^f,\qquad
 \mu(S_R)=\mu(S)-\frac{f\log(T-1)}{T-2},\qquad
 \mu_{\max}(S_R)\le\mu_{\max}(S).
\]
In particular the loss of rank is at most \(f d_S/(T-1)\).
\end{proposition}

\begin{proof}
\proofstep{Show that the residue diamond group acts freely on the fine classes.}
In the split Hermitian model,
let \(K_v^1\) fix the chosen residue row at \(v\mid T\) and
\(K_v^0\) preserve its line. These are subgroups of the
stabilizer of the self-dual overlattice. The scalar on the
row gives \(K_v^0/K_v^1\simeq\mathbf F_T^\times\). It is
surjective because in a basis adapted to the row one can
lift any nonzero scalar by an integral diagonal matrix.
The quotient
\[
                   H=\prod_{v\mid T}\mathbf F_T^\times
\]
acts by right translations on the fine class set, and on
its spaces of values. The other local levels are the same
at both levels. A fixed class for an element of \(H\)
would give a rational stabilizer at the level with local
groups \(K_v^0\), by the equation identifying the translated
adelic representative with the original class. This stabilizer is finite by definiteness and discreteness of
the rational lattice stabilizer.

At a place above \(g_i\) it embeds
in a conjugate of a pro-\(g_i\) group, by the unitriangular residue
condition and the principal congruence filtration.
If \(W\) denotes the value lattice on a representative fibre, the orbit realization reads
\[
  S\simeq\bigoplus_{h\in H}W,\qquad \operatorname{rank} S=|H|\operatorname{rank}W .
\]
 A nontrivial
finite subgroup would contain an element of order \(g_i\). Its
action over \(\mathbb Q\) in dimension \(2fk\) would contain the
cyclotomic polynomial of degree 
\[
  g_i-1>2fk
\]
 in its minimal
polynomial. This is impossible. The inequality holds at all
sufficiently late points since \(g_i\ge(kb)^8\) and \(f\le k\). The class equation then puts the translating
element in the fine level. Thus the \(H\)-action on fine
classes is free.

\proofstep{Transport each integral fibre and its metric along this action.}
Let \(\Lambda_c\) be the value lattice in a basis of the dual
Hermitian lattice at the class \(c\). In the enlarged stabilizer
a diamond \(h\) preserves the underlying lattice. The class
equation identifies it locally at every prime, hence globally.
Let \(U_h\) be transport by the resulting rational isometry
and the level character. The latter is a root of unity. At
finite and infinite places, respectively, we get
\[
 U_h\Lambda_c=\Lambda_{ch},\qquad
 \langle U_h x,U_h y\rangle_{ch}=\langle x,y\rangle_c .
\]
The first equality uses the local unit character and inverse
transport. The second uses the transported metric and the
absolute value one of that character. Distinct class summands
are orthogonal, so transport on their sum is an integral isometry.

Take one class in each \(H\)-orbit, and let \(W\) be the
orthogonal sum of the value lattices on those classes.
Free orbit transport therefore gives an orthogonal integral decomposition
\[
  S\simeq\bigoplus_{h\in H}W.
\]
 On that sum the
diamond operators permute the \(H\)-coordinates by
translations. For \(v\mid T\) write \(N_v\) for the sum
of its diamond operators. The ramified space \(R\) is the
intersection of all the kernels of \(N_v\) in the level-zero
space. Its integral lattice \(S_R=S\cap R\) corresponds to
arrays with sum zero in each one of the \(f\) coordinates.
Put \(n=T-1\). On one factor the sum-zero lattice

\[
  A=\ker(\mathbb Z^n\to\mathbb Z)
\]
 has the basis \(e_i-e_n\) for
\(1\le i<n\). The split decomposition of \(\mathbb Z^n\) into
\(A\) and \(\mathbb Z e_n\), with the sum equal to the projection
on the second factor, works over \(\mathcal O_K\).
Applying it in each coordinate identifies \(S_R\) with
\(W\otimes_{\mathbb Z} A^{\otimes f}\), including its induced
metric, and gives
\begin{equation}\label{eq:ramified-rank}
        d_{\rm ram}:=\operatorname{rank} S_R
           =d_S\left(1-\frac1{T-1}\right)^f .
\end{equation}
This is the dimension of the column after imposing the trace
condition at every place above \(T\).

\proofstep{Compute the determinant of the tensor product of trace-zero factors.}
The Gram matrix of \(A\) in this basis is \(I+J\), where
\(J\) has every entry one. Its determinant is \(n\), since
its eigenvalues are \(n\) on the constants and \(1\) on
their orthogonal complement. The tensor Gram matrix has
size \((n-1)^f\), and its log determinant divided by this
size is \(f\log n/(n-1)\), by the determinant formula for
a tensor product. The integral coordinates above are
free in the \(A\)-factors, so the finite determinant
contribution per rank is the one for \(W\).

At infinity
the squared-norm determinant of the tensor Gram matrix
subtracts \(f\log n/(n-1)\) per rank in the degree. Thus
\begin{equation}\label{eq:ramified-slope}
        \mu(S_R)=\mu(S)-\frac{f\log(T-1)}{T-2},\qquad
        \mu_{\max}(S_R)\le\mu_{\max}(S).
\end{equation}
The inequality uses the induced metric and inclusion of
\(S_R\) in \(S\), so any sublattice used for the maximum
slope has the same degree as a sublattice of \(S\).
For \(T\ge3\) the lost dimension is at most \(f d_S/(T-1)\).
The exact rank and slope formulas apply at every selected sequence with the class and lattice identifications above.
In particular \(T\ge b^{m/2}\ge b\) and
\(f\le k\le x^5\) make both \(f/(T-1)\) and
\(f\log(T-1)/(T-2)\) tend to zero uniformly. For the
second term use that \(\log y/(y-1)\) decreases for \(y>1\).
\end{proof}

\subsection{Weight change at the prime assigned to the sequence}
\label{sec:ramified-weight-change}

At the assigned prime the change of weight alters the integral frames. We compute that change before comparing the two Hecke maps on the trace-zero space.

Let \(w\) be a coefficient-field place over \(T\) at which
the split level comparison is used. It designates one
embedding \(\tau=\tau(w,v)\) for each \(v\mid T\).
Write \(\mathcal O=\mathcal O_{K,w}\) and \(e_w=v_w(T)\),
with \(v_w\) taking value one on a uniformizer. We use
the two coefficient lattices in the actual level-zero
class coordinates. The usual one is \(S_w\). The second,
denoted by \(B_{w,v}\), uses the mixed lattice for the
\(v\)-move in the \(\tau\)-factor and the usual lattice
in every other factor.

\begin{proposition}[Integral change to the mixed column lattice]\label{prop:ramified-mixed-column}
For \(w\mid T\) and the specified split factor at a selected
embedding \(\tau\), the mixed coefficient image is
\(\rho_\lambda(A_c^{-1})\mathcal L_c\) in each class.
With \(b_v=\max(0,-\lambda_k)\), the inclusion and its length from \eqref{eq:column-branch-length} are
\[
 T^{b_v}S_w\subseteq B_{w,v},\qquad \operatorname{length}_{\mathcal O}(B_{w,v}/T^{b_v}S_w)=e_wd_S(s_\lambda/k+b_v).
\]

\end{proposition}

\begin{proof}
\proofstep{The mixed lattice is the marked index-\(T\) neighbor.}

In the indicated split component write \(L_c\) for the
underlying column lattice of a class, so that the
ordinary lattice for the coordinate comparison is
\(L_c^\vee\). The mixed lattice is \(L_c+T L_c^\vee\)
at the indicated place. Its plus factor is the kernel
of the chosen residue row on \(L_c^\vee\), since
\(T L_c^\vee\subseteq L_c\). Thus it is an index-\(T\)
neighbor of \(L_c^\vee\) in the relevant factor.
In a basis whose last
vector has residue value one for that functional, a
basis of the neighbor is given by multiplying that last
vector by \(T\) and keeping the other vectors. This
also proves the determinant valuation for arbitrary
local bases. Denote by \(A_c\) the ordinary-to-mixed
basis matrix of the neighbor, written in ordinary
coordinates at a class \(c\). Then \(A_c\) is integral
and 
\[
  v_w(\det A_c)=e_w.
\]

\proofstep{Apply the divided-root Weyl lattice to the basis change.}

Use the coefficient model of highest weight
\(\lambda_\tau=(\lambda_1,\ldots,\lambda_k)\). Write
\(a=\lambda_k\), \(s_\lambda=\sum_i\lambda_i\), and
\[
       \rho_\lambda(A)=\det(A)^a\,\rho_{\lambda-a\boldsymbol1}(A).
\]
The last representation is polynomial. Its integral model
here is the divided-root Weyl lattice from
Section~\ref{subsec:selected-column-metric}. Unimodular
matrices preserve that lattice by integral Weyl covariance.
For an integral matrix of nonzero determinant over the DVR,
take a Smith factorization into two unimodular matrices
and a diagonal integral matrix. The torus weight spaces of
the polynomial Weyl lattice have nonnegative exponents, so
the diagonal factor also preserves it by inclusion.

In ordinary coordinates at \(c\), write \(\mathcal L_c\)
for the tensor of integral coefficient models at all
embeddings. Write \(\mathcal L_c^{\rm mix}\) for the mixed lattice in these ordinary coordinates. Since the basis change acts only at \(\tau\), its image is
\[
  \mathcal L_c^{\rm mix}=\rho_\lambda(A_c^{-1})\mathcal L_c.
\]
This is an equality of lattices in the rational coefficient
space. The level-zero class realization allows independent
values at the fine classes. The polynomial model is
defined over \(K\), and a basis change through the chosen
embedding of the CM field is a matrix over \(K\).
Thus a polynomial in the indicated image at each class
gives a rational class value.

In any other adapted basis
the transition is unimodular in its local factor, so this
same value satisfies the integral condition in that basis.
Conversely a mixed integral value expressed in a chosen
adapted basis lies in the indicated image. In particular
both coordinate lattices have total rank \(d_S\).

If \(a\ge0\), then \(\rho_\lambda(A_c)\mathcal L_c\subseteq
\mathcal L_c\), giving \(S_w\subseteq B_{w,v}\).
If \(b_v=-a\ge0\), then
\[
 T^{b_v}\det(A_c)^{-b_v}
          =\bigl(T/\det(A_c)\bigr)^{b_v}
\]
is a unit in \(\mathcal O\). It follows that
\(T^{b_v}\rho_\lambda(A_c)\mathcal L_c\subseteq\mathcal L_c\),
and hence \(T^{b_v}S_w\subseteq B_{w,v}\).
\proofstep{Take the determinant in the coefficient representation.}

For an exponent \(b_v\) in either comparison, taking
\(b_v=0\) in the first one, the resulting length is
\begin{equation}\label{eq:column-branch-length}
          \operatorname{length}_{\mathcal O}
                  (B_{w,v}/T^{b_v}S_w)
             = e_w d_S\left(\frac{s_\lambda}{k}+b_v\right).
\end{equation}
Indeed on the irreducible factor of dimension \(d_\tau\),
the determinant of \(\rho_\lambda\) is a power of the
ordinary determinant. This follows by restricting to
the diagonal torus, where its exponent is invariant
under permutations, and observing that the elementary
unipotents act with determinant one. On scalar matrices
the weight is \(s_\lambda\), so the exponent is
\(d_\tau s_\lambda/k\). Tensoring with the other factors
multiplies it by their ranks. Summing the determinant
valuations over classes then gives \eqref{eq:column-branch-length}.
The ramification index comes from \(v_w(T)\).

\end{proof}

\subsection{The integral Hecke index on the trace-zero space}
\label{sec:ramified-hecke-index}

The frame change accounts for the weight of a single term. On the trace-zero space the up and down sums preserve the lattice. Their scalar compositions give the index bound.

\begin{proposition}[The two integral Hecke maps on the trace-zero part]\label{prop:ramified-hecke-maps}
For the indicated split place and coefficient embedding the
up and down sums preserve the integral coordinates as in
\eqref{eq:ramified-hecke-input}. Both preserve the trace-zero space \(R\). On this space the two compositions satisfy
\[
 (D_vU_v)|_R=(U_vD_v)|_R=T^{k-1}1_R.
\]

\end{proposition}

\begin{proof}
\proofstep{Express the two transition maps in the mixed frames.}

In the local comparison let \(U_v\) and \(D_v\) be the
finite sums for the up and down moves. Use the chosen split
component over \(v\), and write \(\mathcal O_v\) for its
integer ring, with uniformizer \(T\). A local state has a
lattice \(L\) and an index-\(T\) overlattice \(D\), the
relevant component of \(L^\vee\), with a marking
\(\lambda:D/TD\to\mathbf F_T\) whose kernel is \(L/TD\).
The pairing recovers the lattices in the opposite split
component from \(L,D\). Choose \(z_0\in D\) with
\(\lambda(z_0)=1\) and let \(t=Tz_0\bmod TL\) in \(L/TL\).
This vector is independent of the choice of \(z_0\), and
is nonzero since \(z_0\notin L\).

In the chosen split bases the up move with \(z\bmod TD\)
and \(\lambda(z)=1\) has lattices and marking
\[
           D\ \subset\ D+\mathcal O_v T^{-1}z,\qquad
                        z'\longmapsto\lambda(Tz').
\]
For the down move choose \(g\in(L/TL)^*\) with \(g(t)=1\). Its lattice inclusion and residue marking are
\[
  L_g=\{x\in L:g(x\bmod TL)=0\}\subset L,\qquad L/L_g\xrightarrow{\ g\ }{\cal O}_v/T.
\]
These descriptions follow from the up rule of adjoining
the line generated by \(T^{-1}z\) to the ordinary dual
lattice, and the down rule of taking the marked residue
kernel on the mixed lattice. At other finite places the
move is trivial. In class coordinates these equations are
read after the rational isometry that transports the new
class to the moving lattice. The two endpoint descriptions
also show the lattice integrality.

For an up move the
transported mixed lattice at the new class is \(D\),
the ordinary dual lattice of the starting class.
Indeed, in the enlarged lattice the transported marking has kernel \(D\). On generators its values are
\[
  \lambda(TD)=0\pmod T,\qquad \lambda\!\left(T\cdot T^{-1}z\right)=1\pmod T.
\]
 For a
down move the transported ordinary dual lattice is \(L\),
the mixed lattice of the starting class. At the other
embeddings the ordinary basis lattices agree under the
transport. Thus the polynomial coefficient maps carry
integral values to integral values, by changing adapted
bases as in Section~\ref{sec:ramified-weight-change}.
The value in the transported frame is multiplied only by
the level character, a unit in \(\mathcal O\). The two integral transport maps therefore satisfy
\[
  U_v B_{w,v}\subseteq S_w,\qquad D_vS_w\subseteq B_{w,v}.
\]

\proofstep{Count the end states of the composite.}

To compute the composite \(D_v U_v\), first choose a down
label \(g\), then an up vector \(y\in L/TL\) with \(g(y)=1\).
At \(v\) the end state is the pair
\[
         L\ \subset\ L+\mathcal O_v T^{-1}\widetilde y,
                \qquad z'\longmapsto g(Tz'\bmod TL).
\]
Here \(\widetilde y\) is any lift into \(L\). The marking is
zero on \(L\) and sends \(T^{-1}\widetilde y\) to one, so
this state and marking depend only on \(y\). The support
of the two moves is \(v\), and the data of a split lattice
and its residue marking determine the coset of the local
level.

The local level character is trivial at this prime.
Consequently, when evaluating at the starting frame, pairs
with the same \(y\) give the same automorphic value
\(F_y\). In particular \(F_t\) is the original value, and
\(F_{a t}\) is its diamond translate by the scalar \(a^{-1}\),
since the marking on the original \(D\) now sends \(az_0\)
to one. Inversion permutes the nonzero scalars.

For nonzero \(y\), use the evaluation map

\[
  g\mapsto(g(t),g(y))
\]
 on the dual \(k\)-dimensional vector space.
The fibre of \((1,1)\) has size
\[
 \#\{g:g(t)=g(y)=1\}=
 \begin{cases}
 T^{k-1}, &y=t,\\
 0, &y=at,\ a\ne0,1,\\
 T^{k-2}, &t,y\text{ linearly independent}.
 \end{cases}
\]
The evaluation has rank one in the first two cases. In the second
case its image is the line \((s,as)\), which omits \((1,1)\).
In the third case its rank is two.
For the corresponding sum \(E_v U_v\) allowing all
nonzero \(g\) in the first move, every nonzero \(y\)
has \(T^{k-1}\) choices. Thus for \(P_v=T^{-(k-1)}E_vU_v\)
one obtains
\[
       D_v U_v=T^{k-1}1-T^{k-2}N_v+T^{k-2}P_v .
\]
The normalized trace \(P_v\) sums the states \(F_y\)
over the nonzero vectors in \(L/TL\). A diamond on
these states rescales their markings, which permutes
the nonzero \(y\), so \(P_v\langle a\rangle_v=P_v\).
\proofstep{Pass to the trace-zero part using the diamond actions.}

Also \(U_v,D_v\) commute with a diamond at \(v\).
Indeed \(\lambda\mapsto a\lambda\) is accompanied by
\(z\mapsto a^{-1}z\) in the up sum and by \(g\mapsto ag\)
in the down sum. The endpoint lattices agree and the
endpoint markings are multiplied by \(a\) in both cases.
Diamonds at the other places commute by disjoint support.
Summing the trace identity over \(a\) gives

\[
  (T-1)P_v=P_vN_v.
\]
 On \(R\) both \(N_v\) and \(P_v\)
vanish, and the diamond commutations show that both
Hecke operators preserve \(R\). The composite \(D_vU_v\)
on \(R\) is \(T^{k-1}1_R\). In finite dimension this
makes \(U_v|_R\) invertible and gives \(U_vD_v|_R=T^{k-1}1_R\).
We have proved
\begin{equation}\label{eq:ramified-hecke-input}
 \begin{gathered}
 U_v B_{w,v}\subseteq S_w,\qquad D_v S_w\subseteq B_{w,v},\\
 U_v R,D_v R\subseteq R,\qquad U_v D_v|_R=T^{k-1}1_R .
 \end{gathered}
\end{equation}

\end{proof}

\begin{proposition}[Colength of the selected higher-level image]\label{prop:ramified-selected-hecke-colength}
With \(M=S_w\cap R\), the normalized up map \(G_v\) is injective on \(M\). The index bound \eqref{eq:ramified-hecke-index} is
\[
 \operatorname{length}_{\mathcal O}(M/G_vM)\le e_wd_S(s_\lambda/k+b_v+k-1).
\]

The colength for the assigned product of the \(G_v\)'s to the
\(m\)-th power is at most the sum of the individual bounds times
\(m\), and is zero for an empty product. The weighted logarithmic
degree over \(T\) uses \(\sum_{w\mid T} e_w\omega_w=\log T\).
\end{proposition}

\begin{proof}
\proofstep{Compare the two integral images inside the trace-zero space.}

To see which index is needed, set
\(M=S_w\cap R\) and \(B_R=B_{w,v}\cap R\). For
\(b_v=\max(0,-a)\), put \(G_v=T^{b_v}U_v|_R\).
The mixed lattice inclusion gives \(T^{b_v}M\subseteq B_R\).
The two Hecke lattice maps and their composition give
\[
         T^{k-1}M\subseteq U_v B_R\subseteq M,\qquad
                           G_v M\subseteq U_v B_R .
\]
Consequently the colength of \(G_v M\) in \(M\) is the
sum of the colength of \(U_v B_R\) in \(M\) and the
length of \(B_R/T^{b_v}M\), using injectivity of \(U_v\)
on \(R\).

The quotient \(B_R/T^{b_v}M\) embeds into \(B_{w,v}/T^{b_v}S_w\), since
\(T^{b_v}S_w\cap R=T^{b_v}M\).
The first colength is at most
\(e_w(k-1)\operatorname{rank} M\), since the quotient
is a quotient of \(M/T^{k-1}M\). Formula
\eqref{eq:column-branch-length} therefore gives
\begin{equation}\label{eq:ramified-hecke-index}
 \operatorname{length}_{\mathcal O}(M/G_vM)
           \le e_w d_S\left(\frac{s_\lambda}{k}+b_v+k-1\right).
\end{equation}
This accounts for the ordinary and ramified lattice
comparisons in the same coefficient space.

\proofstep{Add the colengths of the injective factors.}

In the column restriction at \(w\mid T\), the selected
places \(v\) are those assigned the top twists.
The unshifted column lattice at \(w\) is the image of \(M\)
by the product of their \(G_v\), raised to the \(m\)-th power.
Each factor is injective with \(G_vM\subseteq M\).

For
any two such factors \(G,G'\) the chain
\(GG'M\subseteq GM\subseteq M\) gives the sum
of their colengths, since \(G\) identifies
\(M/G'M\) with \(GM/GG'M\). Induction and
\eqref{eq:ramified-hecke-index} bound the local colength by
\[
          m e_w d_S\sum_{v}
                 \left(\frac{s_{\lambda_{\tau(w,v)}}}{k}
                                      +b_v+k-1\right)
\]
with the sum over those selected \(v\). The product is the
identity in the empty case, in particular at the places
where the split level comparison is unused. The estimate
has the same coefficients at each place. At a top twist
let \(a_{\max}\) bound the first weight coordinate and
let \(W_\tau\) bound the absolute values of the other
coordinates. Then \(b_v\le W_\tau\) and

\[
  s_{\lambda_\tau}\le a_{\max}+(k-1)W_\tau.
\]
The term in parentheses is bounded by
\(a_{\max}/k+2W_\tau+k-1\). These are the weight
bounds for the given top twists, with \(\tau(w,v)\)
distinct as \(v\) varies at each \(w\).
\proofstep{Sum over coefficient places without a degree loss.}

Passing to the logarithmic degree uses the weighted
place identity \(\sum_{w\mid T}e_w\omega_w=\log T\),
so this step introduces no coefficient-field degree
factor. For example the difference of slopes from the
indices over \(T\) is bounded by
\[
 \frac{2m d_S\log T}{d_{\rm ram}}
     \sum_{\tau\ {\rm top}}
                  \left(a_{\max}/k+2W_\tau+k-1\right).
\]
The ratio \(d_S/d_{\rm ram}\) is bounded by
\((1-f/(T-1))^{-1}\) once \(f/(T-1)<1\), hence
by \(2\) above a common scale cutoff.


\label{cor:ramified-column-1}
\end{proof}
